\chapter{Bande de particule, anti-section et virtualité}
\label{chap:banda-particula-virtual}
\index{particule!bande de parcours}
\index{antiparticule!anti-section orientée}
\index{particule virtuelle}

La persistance d’un parcours produit une particule avant l’évaluation de sa masse. Cette priorité permet de distinguer trois objets :

\begin{enumerate}
\item la \emph{section orientée}, qui conserve la direction du parcours ;
\item la \emph{bande}, qui réunit une section et sa conjuguée ;
\item la \emph{persistance}, qui décide si la signature atteint ou non une branche globale.
\end{enumerate}

\begin{definition}[Mots dodécaphasiques et conjugaison matérielle]
\label{def:palabras-cm}
Soit
\[
 \mathscr T_{12}^{\rm word}:=\{-1,0,1\}^{12}.
\]
Pour \(\mathbf t=(t_1,\ldots,t_{12})\), nous écrivons
\[
 \operatorname{rev}(\mathbf t)
 =(t_{12},\ldots,t_1),
 \qquad
 C_M(\mathbf t):=-\operatorname{rev}(\mathbf t).
\]
L’inversion \(C_M\) change simultanément le sens de lecture et l’orientation tritique.
\end{definition}

\begin{theorem}[Involution matérielle et secteur fixe]
\label{thm:involucion-cm-censo} L’application \(C_M\) est une involution de \(\mathscr T_{12}^{\rm word}\) et
\[
 \left|\operatorname{Fix}(C_M)\right|=3^6=729.
\]
\end{theorem}

\begin{proof}
Le renversement et le changement de signe commutent et ont tous deux pour carré l’identité ; par conséquent, \(C_M^2=I\). Un mot est fixe si et seulement si
\[
 t_i=-t_{13-i},\qquad 1\leq i\leq6.
\]
Les six premières coordonnées sont libres et les six dernières sont déterminées par elles. Chaque coordonnée libre admet trois valeurs, de sorte que le secteur fixe a pour cardinal \(3^6\).
\end{proof}

Cette conjugaison finie ne s’obtient pas en changeant uniquement le signe visible : elle inverse aussi l’orientation de l’enregistrement.

\section{La bande de parcours}

Soit \(\mathscr T_{12}^{\rm word}\) l’espace de mots de la \cref{def:palabras-cm}. Nous ajoutons une orientation de section \(\chi\in\{-1,+1\}\) et définissons
\[
(\mathbf t,\chi)\sim(C_M\mathbf t,-\chi).
\]

\begin{definition}[Bande dodécaphasique de parcours]
\label{def:banda-dodecafasica-ruta} La bande de la section orientée \((\mathbf t,\chi)\) est la classe
\[
[(\mathbf t,\chi)]_M
=\{(\mathbf t,\chi),(C_M\mathbf t,-\chi)\},
\]
et l’espace des bandes est
\[
\mathcal B_{12}
:=
\bigl(\mathscr T_{12}^{\rm word}\times\{-1,+1\}\bigr)/{\sim}.
\]
\end{definition}

\begin{proposition}[Liberté de la conjugaison orientée]
\label{prop:accion-cm-orientada-libre}
L’action
\[
(\mathbf t,\chi)\longmapsto(C_M\mathbf t,-\chi)
\]
est libre. Par conséquent,
\[
|\mathcal B_{12}|=3^{12}.
\]
\end{proposition}

\begin{proof}
Un point fixe exigerait \(\chi=-\chi\), ce qui est impossible pour \(\chi\in\{-1,+1\}\). Toutes les orbites ont deux éléments et le domaine possède \(2\cdot3^{12}\).
\end{proof}

La bande n’efface pas l’orientation : elle la conserve comme revêtement double. La particule et l’antiparticule sont les deux sections orientées de la même bande.

\section{Masse paire et charge impaire}

Soit \(E_0:\mathscr T_{12}^{\rm word}\to\mathbb R\) un caractère additif de parcours. Ses composantes paire et impaire par rapport à \(C_M\) sont
\[
E_+(\mathbf t)
=\frac12\bigl(E_0(\mathbf t)+E_0(C_M\mathbf t)\bigr),
\]
\[
E_-(\mathbf t)
=\frac12\bigl(E_0(\mathbf t)-E_0(C_M\mathbf t)\bigr).
\]
Alors
\[
E_+(C_M\mathbf t)=E_+(\mathbf t),
\qquad
E_-(C_M\mathbf t)=-E_-(\mathbf t).
\]

\begin{definition}[Caractère de bande]
\label{def:caracter-banda} Pour une section orientée, nous définissons
\[
E_{\rm band}(\mathbf t,\chi)
=E_+(\mathbf t)+\chi E_-(\mathbf t),
\qquad
m(\mathbf t,\chi)
=m_\circ\exp\bigl(E_{\rm band}(\mathbf t,\chi)\bigr),
\]
où \(m_\circ\) est l’unité interne de comparaison de la famille.
\end{definition}

\begin{theorem}[Principe de conjugaison de masse]
\label{thm:principio-conjugacion-masa-banda} Le caractère et la masse passent au quotient des bandes :
\[
E_{\rm band}(C_M\mathbf t,-\chi)
=E_{\rm band}(\mathbf t,\chi),
\]
\[
\boxed{
m(C_M\mathbf t,-\chi)=m(\mathbf t,\chi).}
\]
Les invariants pairs appartiennent à la bande ; les invariants impairs distinguent l’orientation de ses deux sections.
\end{theorem}

\begin{proof}
D’après les parités précédentes,
\[
\begin{aligned}
E_{\rm band}(C_M\mathbf t,-\chi)
&=E_+(C_M\mathbf t)-\chi E_-(C_M\mathbf t)\\
&=E_+(\mathbf t)+\chi E_-(\mathbf t).
\end{aligned}
\]
L’égalité des masses s’obtient en appliquant l’exponentielle.
\end{proof}

En particulier, une forme linéaire à poids symétriques est impaire, tandis qu’une forme quadratique compatible avec le renversement est paire. La charge et l’orientation changent ; la masse, construite sur la bande, demeure.

\begin{corollary}[Électron et positron]
\label{cor:electron-positron-banda} Dans l’équation électronique de bande, la paire \((n,\chi)=(-22,+1)\) et sa conjuguée \((+22,-1)\) produisent le même terme \(\chi n=-22\). L’électron et le positron possèdent donc la même masse et des orientations de charge opposées.
\end{corollary}

\section{Clôture finie de l’atlas sous anti-section}

Dans les quatre-vingt-un parcours CT--108 apparaissent sept formes tritiques. Le dénombrement exhaustif distingue trois opérations :
\[
\mathbf t\longmapsto-\mathbf t,\qquad
\mathbf t\longmapsto\operatorname{rev}(\mathbf t),\qquad
\mathbf t\longmapsto-\operatorname{rev}(\mathbf t).
\]

\begin{proposition}[Clôture des formes CT--108]
\label{prop:cierre-antiseccion-ct108} L’ensemble des formes tritiques de CT--108 n’est stable ni sous le seul changement de signe ni sous le seul renversement, et il est stable sous \(C_M=-\operatorname{rev}\). Ses orbites possèdent les multiplicités
\[
\boxed{81=54+9+9+9.}
\]
L’orbite de multiplicité \(54\) est autoconjuguée ; les trois autres échangent des paires de formes.
\end{proposition}

\begin{proof}
La forme fixe est
\[
\texttt{+++---+++---}
\]
et possède la multiplicité \(54\). Les trois paires
\[
\begin{array}{c}
\texttt{+++--0+++--0}\ /\ \texttt{0++---0++---},\\
\texttt{+++-0-+++-0-}\ /\ \texttt{+0+---+0+---},\\
\texttt{+++0--+++0--}\ /\ \texttt{++0---++0---}
\end{array}
\]
possèdent chacune la multiplicité totale \(9\). L’application \(-\operatorname{rev}\) fixe la première et permute les membres de chaque paire. Le changement de signe et le renversement, pris séparément, produisent des mots extérieurs à cette liste. La somme des multiplicités est \(54+9+9+9=81\).
\end{proof}

Le résultat n’identifie pas le secteur \(54\) à un autre objet de même cardinal. Il démontre quelle conjugaison respecte la taxonomie concrète des parcours et pourquoi particule et antiparticule partagent une bande.

\section{Persistance finie et particule virtuelle}

Soit
\[
\cdots\xrightarrow{\rho_{M+1,M}}S_M
\xrightarrow{\rho_{M,M-1}}S_{M-1}\xrightarrow{}\cdots
\]
le système des états admissibles aux profondeurs finies. Pour \(s_N\in S_N\), nous définissons l’ensemble des prolongements
\[
\mathcal P_M(s_N)
:=
\{s_M\in S_M:\rho_{M,N}(s_M)=s_N\},
\qquad M\geq N.
\]
Sa profondeur de vie est
\[
\mathcal L(s_N)
:=\sup\{M\geq N:\mathcal P_M(s_N)\neq\varnothing\}
\in\mathbb N\cup\{\infty\}.
\]

\begin{definition}[Persistance et virtualité]
\label{def:persistencia-virtualidad} Un préfixe \(s_N\) est \emph{persistant} s’il appartient à une section globale \((s_M)_{M\geq N}\) compatible. Il est \emph{virtuel fini} si \(\mathcal L(s_N)<\infty\), s’il existe un prolongement jusqu’à \(\mathcal L(s_N)\) et s’il n’en existe aucun à la profondeur suivante.
\end{definition}

\begin{theorem}[Critère de persistance]
\label{thm:criterio-persistencia-virtual} Supposons que chaque état possède un nombre fini de prolongements immédiats et que les restrictions soient cohérentes. Alors :
\[
\mathcal L(s_N)=\infty
\quad\Longleftrightarrow\quad
s_N\text{ appartient à une section globale}.
\]
Si \(\mathcal L(s_N)=L<\infty\), le préfixe transporte le parcours, la mémoire, l’enregistrement et la signature jusqu’à \(L\), mais ne définit pas une particule persistante.
\end{theorem}

\begin{proof}
Une section globale fournit un prolongement à toute profondeur, de sorte que \(\mathcal L=\infty\). Réciproquement, les prolongements de \(s_N\) forment un arbre enraciné, infini et finiment ramifié. D’après le lemme de Kőnig, il contient une branche infinie, dont les sommets sont compatibles par construction et définissent une section globale.

Si \(\mathcal L=L<\infty\), il existe un état à la profondeur \(L\) et aucun à \(L+1\). Les lecteurs définis sur les préfixes peuvent l’évaluer jusqu’à \(L\) ; l’absence de prolongement empêche de transformer cette signature finie en une section de la limite inverse.
\end{proof}

La particule virtuelle possède ainsi une définition interne. Elle n’est ni une particule ``moins réelle'' ni une brève violation d’une loi de conservation. C’est un parcours dont la mémoire agit pendant une profondeur finie et dont la généalogie s’éteint avant de se stabiliser. La confrontation avec les propagateurs et les états intermédiaires de la physique conventionnelle intervient après l’identification du lecteur qui représente ces profondeurs.

\section{Chaîne parcours–enregistrement–signature–caractère–tour de la réalisation matérielle}

Les résultats de ce chapitre fixent la direction complète :
\[
\boxed{
\text{extension}\to
\text{route}\to
\text{registre}\to
\text{signature}\to
\text{bande orientée}\to
\text{particule}\to
\text{masse}.}
\]
L’antiparticule est la seconde section de la bande. La virtualité est une persistance finie. Aucune de ces notions n’est définie à partir d’une masse observée.
